\section*{Capítulo \ref{ch:b2:retrosynthesis} --- Retrosíntesis y estrategia en varias etapas}

\begin{solution}{exo:b2:retrosynthesis:1}
\begin{itemize}
  \item 2-Fenilpropan-2-ol: propanona + \ce{PhMgBr}, o acetofenona + \ce{CH3MgBr}.
  \item Pentan-3-ol: propanal + \ce{CH3CH2MgBr}.
  \item 1-Ciclohexiletanol: etanal + bromuro de ciclohexilmagnesio, o ciclohexanocarbaldehído +
        \ce{CH3MgBr}.
  \item 2-Metilbutan-2-ol: propanona + \ce{CH3CH2MgBr}, o butanona + \ce{CH3MgBr}.
  \item 1-Feniletanol: benzaldehído + \ce{CH3MgBr}, o etanal + \ce{PhMgBr}.
\end{itemize}
\end{solution}

\begin{solution}{exo:b2:retrosynthesis:2}
\omterm{def:b2:retrosynthesis:disconnection}{Sintón} aceptor $\mathrm{Ph\overset{+}{C}H{-}OH}$, equivalente el benzaldehído; \omterm{def:b2:retrosynthesis:disconnection}{sintón} dador \ce{CH3-},
equivalente el bromuro de metilmagnesio (o el metil-litio). (El otro corte, con \ce{Ph-} como dador, usa
\ce{PhMgBr} y etanal.)
\end{solution}

\begin{solution}{exo:b2:retrosynthesis:3}
$0.90 \times 0.85 \times 0.80 \times 0.95 \times 0.70 = 0.407$: un 41~\% aproximadamente.
\end{solution}

\begin{solution}{exo:b2:retrosynthesis:4}
El \ce{LiAlH4}, necesario para reducir el éster al alcohol, reduciría también la cetona. Protéjase la
cetona como acetal cíclico (etano-1,2-diol, catalizador ácido, eliminación del agua); redúzcase el éster con
\ce{LiAlH4}; elimínese el acetal con ácido acuoso diluido: 5-hidroxipentan-2-ona.
\end{solution}

\begin{solution}{exo:b2:retrosynthesis:5}
IGF: butano-1,3-diol $\Rightarrow$ 3-hidroxibutanal (reducción del aldehído, \ce{NaBH4}). \omterm{def:b2:retrosynthesis:disconnection}{Desconexión} de dos grupos
de este $\beta$-hidroxialdehído en el enlace $\alpha$--$\beta$: dos moléculas de etanal
(\omterm{def:b2:enolates-aldol:aldol}{aldólica}). Hacia delante: etanal, base diluida, en frío; y luego \ce{NaBH4}.
\end{solution}

\begin{solution}{exo:b2:retrosynthesis:6}
El anillo de ciclohexeno se corta por los dos enlaces \ce{C-C} formados en una \omterm{def:b2:frontier-orbitals:diels-alder}{reacción de Diels--Alder}, los
alílicos respecto al doble enlace del anillo del lado del sustituyente: buta-1,3-dieno y but-3-en-2-ona, con el
grupo acetilo en el \omterm{def:b2:frontier-orbitals:diels-alder}{dienófilo}.
\end{solution}

\begin{solution}{exo:b2:retrosynthesis:7}
Numerando C1 (\ce{C=O}), C2=C3, y luego C4 con los dos metilos
(\cref{met:b2:conjugate-additions:robinson-disconnection}): el corte aldólico C2--C3 deja un aldehído
en C3 y una metilcetona; el corte de Michael C4--C5 deja 2-metilpropanal (C3 su carbonilo, C4 su
carbono $\alpha$) y but-3-en-2-ona.
\end{solution}

\begin{solution}{exo:b2:retrosynthesis:8}
Una etapa: benceno con cloruro de butanoílo y \ce{AlCl3} (\omterm{def:b2:aromatic-substitution:friedel-crafts}{acilación de Friedel--Crafts}, sin transposición,
acilación única). Dos etapas: benzaldehído con bromuro de propilmagnesio, y luego oxidación del alcohol
secundario. La ruta de Friedel--Crafts es más corta y sus reactivos más baratos, pero consume un equivalente completo de
cloruro de aluminio, que acaba en los residuos acuosos.
\end{solution}

\begin{solution}{exo:b2:retrosynthesis:9}
El clorocromato de piridinio en diclorometano, la oxidación de Swern o el peryodinano de Dess--Martin
dan hexanal; el dicromato acuoso acidulado o el permanganato dan ácido hexanoico, porque en agua el
aldehído forma su hidrato, que se oxida de nuevo (\cref{prop:b2:retrosynthesis:mild-oxidation}).
\end{solution}

\begin{solution}{exo:b2:retrosynthesis:10}
La amina reacciona primero, sin proteger, si el alcohol está protegido: protéjase el alcohol como éter sililado
(sobrevive a la acilación); acílese la amina; elimínese el grupo sililo con fluoruro; acílese el
alcohol. Si hubiera que mantener libre la amina mientras reacciona el alcohol, se protegería como
carbamato de \textit{terc}-butoxicarbonilo (eliminado con ácido), ortogonal al éter sililado (eliminado con
fluoruro): cualquiera de los dos puede liberarse sin tocar el otro.
\end{solution}

\begin{solution}{exo:b2:retrosynthesis:11}
\ce{CH3COCH2CH2COCH3} es un 1,4-dicarbonilo: córtese el enlace \ce{C-C} central. El dador es el \omterm{def:b2:enolates-aldol:enolate}{enolato} del
3-oxobutanoato de etilo (etóxido de sodio); el aceptor, un carbono $\alpha$ de una cetona, lo proporciona la
1-bromopropan-2-ona, \ce{BrCH2COCH3} (sustitución bimolecular del bromuro). Después, la hidrólisis acuosa del
éster y el calentamiento eliminan \ce{CO2}: hexano-2,5-diona.
\end{solution}

\begin{solution}{exo:b2:retrosynthesis:12}
Retrosíntesis: la metilcetona con un \ce{CH2} contiguo es el producto de la síntesis
acetilacética; córtese el enlace entre C3 y C4: el \omterm{def:b2:enolates-aldol:enolate}{enolato} del 3-oxobutanoato de etilo y el haluro alílico
1-bromo-3-metilbut-2-eno, \ce{(CH3)2C=CHCH2Br}. Hacia delante: etóxido de sodio en etanol, y luego el haluro;
después base acuosa, acidificación y calentamiento, que hidrolizan el éster y eliminan \ce{CO2}
(\cref{prop:b2:enolates-aldol:decarboxylation}): 6-metilhept-5-en-2-ona.
\end{solution}

\begin{solution}{pb:b2:retrosynthesis:1}
\textbf{1.} \ce{HO-C6H4-CH2CH2COCH3}, con el \ce{OH} en para respecto a la cadena: un fenol y una cetona.
\textbf{2.} La hidrogenación de un \ce{C=C}: \ce{ArCH2CH2CO} $\Rightarrow$ \ce{ArCH=CHCO}, una
cetona $\alpha,\beta$-insaturada.
\textbf{3.} (\textit{E})-4-(4-hidroxifenil)but-3-en-2-ona, \ce{HOC6H4CH=CHCOCH3}.
\textbf{4.} Un carbonilo $\alpha,\beta$-insaturado: una \omterm{def:b2:enolates-aldol:aldol}{condensación aldólica}, con el corte en el \ce{C=C}.
\textbf{5.} El 4-hidroxibenzaldehído y la propanona.
\textbf{6.} El aldehído no tiene hidrógeno $\alpha$ y solo puede ser el electrófilo; la propanona, en
exceso, es la única fuente de \omterm{def:b2:enolates-aldol:enolate}{enolato} y reacciona con el aldehído, más reactivo, y no consigo misma.
\textbf{7.} (1) 4-hidroxibenzaldehído, exceso de propanona, hidróxido de sodio acuoso; acidificación. (2)
\ce{H2} sobre paladio sobre carbono, a temperatura ambiente, hasta que se absorbe un equivalente.
\textbf{8.} El grupo \ce{CHO} estabiliza el fenóxido, de modo que este fenol es al menos tan ácido como el fenol
($\mathrm pK_a$ 9.99). A $\mathrm{pH} > 13$, $[\ce{ArO-}]/[\ce{ArOH}] > 10^{13-10} = 10^3$: el
fenóxido.
\textbf{9.} El \ce{O-} cede densidad electrónica a través del anillo al carbono carbonílico (un
dador por resonancia en para respecto a él): el aldehído es menos electrófilo, y la condensación, más lenta.
\textbf{10.} Un éter bencílico: la hidrogenación sobre paladio que reduce el \ce{C=C} rompe también
el éter bencílico (hidrogenólisis), liberando el fenol en la misma etapa.
\textbf{11.} Tres: bencilación, \omterm{def:b2:enolates-aldol:aldol}{condensación aldólica}, hidrogenación con desprotección.
\textbf{12.} No: en condiciones suaves (temperatura ambiente, paladio) el anillo \omterm{def:b2:huckel:aromatic}{aromático} no se reduce.
\textbf{13.} No proteger: el fenóxido solo frena la condensación, cosa que compensan un tiempo más largo o un ligero
calentamiento; ahorrar una etapa vale más que la ganancia en velocidad.
\textbf{14.} El 4-yodofenol y la but-3-en-2-ona.
\textbf{15.} \ce{HOC6H4I + CH2=CHCOCH3 + Et3N -> HOC6H4CH=CHCOCH3 + Et3NH+ + I-}, catalizador de paladio.
\textbf{16.} En el \ce{CH2} terminal, el carbono $\beta$: el grupo arilo se inserta en el extremo menos impedido,
y la eliminación de $\beta$-hidruro da después la (\textit{E})-enona conjugada.
\textbf{17.} \omterm{def:b2:catalytic-cycles:oxidative-addition}{Adición oxidante} del yoduro de arilo al \ce{Pd(0)}, coordinación e inserción del
alqueno, eliminación de $\beta$-hidruro y regeneración del \ce{Pd(0)} por la base, que capta el \ce{HI}
(\cref{ch:b2:catalytic-cycles}).
\textbf{18.} $162.188/(220.005 + 70.091 + 101.193) = 162.188/391.289 \approx 41.4~\%$.
\textbf{19.} \ce{C7H6O2 + C3H6O -> C10H10O2 + H2O}: $162.188/(122.123 + 58.080) = 162.188/180.203
\approx 90.0~\%$.
\textbf{20.} 100~\%: todos los átomos del \ce{H2} y de la enona acaban en el producto.
\textbf{21.} $164.204/(122.123 + 58.080 + 2.016) = 164.204/182.219 \approx 90.1~\%$.
\textbf{22.} El \ce{C=C}: conjugado y no impedido, se une al paladio y se hidrogena mucho más deprisa
que el \ce{C=O} de la cetona; detenerse tras un equivalente de hidrógeno conserva la cetona.
\textbf{23.} Ruta \omterm{def:b2:enolates-aldol:aldol}{aldólica}: aldehído barato, propanona e hidróxido de sodio, con agua como subproducto. Ruta de
Heck: un yoduro de arilo, un catalizador de paladio y la but-3-en-2-ona, muy tóxica.
\textbf{24.} La ruta \omterm{def:b2:enolates-aldol:aldol}{aldólica}: dos etapas, reactivos baratos y más seguros, economía atómica superior al 90~\%.
\textbf{25.} Lo multiplica por 0.90: $0.782 \times 0.90 \approx 70~\%$.
\textbf{26.} $0.85 \times 0.92 = \boldsymbol{\approx 78~\%}$ global.
\end{solution}
